% Build from the repository root with `make paper`.
\documentclass[10pt]{article}
% The text width of the ACM paper (acmsmall), so a display that overflows there overflows here.
\usepackage[textwidth=395.8pt,textheight=9in,hcentering]{geometry}
\usepackage{amsmath,amssymb,amsthm,stmaryrd,mathpartir}
\usepackage{array}
\usepackage[colorlinks=true,linkcolor=blue]{hyperref}
\newtheorem{theorem}{Theorem}
\newtheorem{definition}[theorem]{Definition}
\pagestyle{empty}

% Extracted fragments carry declaration links; the displays show only their math.
\newenvironment{paperfragment}[1][]{\ignorespaces}{\ignorespacesafterend}
\newcommand{\lean}[1]{}
\newcommand{\uses}[1]{}
\newcommand{\leanok}{}
\newcommand{\mathlibok}{}
\newcommand{\discussion}[1]{}
\newcommand{\notready}{\PackageError{semantics}{Unfinished blueprint node}{Complete the Lean declaration before including it.}}
\newcommand{\kw}[1]{\mathsf{#1}}
\newcommand{\eval}[1]{\llbracket #1\rrbracket_{\sigma}}
\newcommand{\ir}[3]{\inferrule[\textsc{#1}]{#2}{#3}}
\newcommand{\syntaxline}[2]{%
  \noindent\makebox[5em][l]{\textsf{#1}}\mbox{$#2$}\par}
% A block of semantics displays: a one-item trivlist in footnote size, so the class's own
% list spacing frames it and the font supplies the display skips. The same definition goes
% to the paper through generated/macros.tex.
\newenvironment{semanticsblock}[1][\footnotesize]{%
  \begin{trivlist}#1
  \def\MathparLineskip{\lineskiplimit=0.6em\lineskip=0.6em plus 0.2em}%
  \item[]\ignorespaces}{\end{trivlist}}
% The named blocks scripts/paper_sync.py cuts the paper into: a display block and a prose
% block. The name is the generated file's name.
\newenvironment{semantics}[2][\footnotesize]{\begin{semanticsblock}[#1]}{\end{semanticsblock}}
\newenvironment{semanticstext}[1]{\par}{\par\medskip}

% A missing annotation is a build error, not an empty display.
\makeatletter
\newcommand{\inputleannode}[1]{%
  \@ifundefined{leannode@#1}%
    {\PackageError{semantics}{Missing Lean node #1}{Check the label and extraction module list.}}%
    {\csname leannode@#1\endcsname}}
\makeatother
\input{.lake/build/blueprint/library/EgglogPaper}

\begin{document}
\begin{semantics}[\small]{syntax}
  \(x\in\mathit{Var}\qquad f\in\mathit{Fn}\)\par
  \noindent\begin{minipage}[t]{.6\linewidth}
    \inputleannode{def:expressions}
    \inputleannode{def:action-syntax}
    \inputleannode{def:commands}
  \end{minipage}\hfill
  \begin{minipage}[t]{.4\linewidth}
    \inputleannode{def:patterns}
    \inputleannode{def:queries}
    \inputleannode{def:rules}
    \inputleannode{def:programs}
  \end{minipage}
\end{semantics}

\begin{semantics}{egraph}
  {\centering\inputleannode{def:terms}\qquad\inputleannode{def:state}\par}
  \inputleannode{def:congruence}
  \inputleannode{def:stored-terms}
  \inputleannode{def:live}
  \inputleannode{def:equiv}
\end{semantics}

\begin{semanticstext}{order}
  \inputleannode{def:order}
\end{semanticstext}

\begin{semantics}{matching}
  \inputleannode{def:representatives}
  \inputleannode{def:query}
\end{semantics}

\begin{semantics}{evaluation}
  \inputleannode{def:evaluation}
  \inputleannode{def:actions}
\end{semantics}

\begin{semantics}{execution}
  \inputleannode{def:round}
  \inputleannode{def:saturation}
  \inputleannode{def:behavior}
\end{semantics}

\begin{semanticstext}{wellformed}
  \inputleannode{def:wellformed}
  \inputleannode{def:fragment}
  \inputleannode{thm:congruence}
\end{semanticstext}

% The congruence pass as inference rules; the paper's congruence section inputs these.
\begin{semantics}{compile-query}
  \inputleannode{def:compile-query}
\end{semantics}

\begin{semantics}{compile-head}
  \inputleannode{def:compile-head}
\end{semantics}

\begin{semantics}{maintenance}
  \inputleannode{def:maintenance}
\end{semantics}

\begin{semantics}{compile-schedule}
  \inputleannode{def:compile-schedule}
\end{semantics}

% The optimizations of the congruence pass: canonical node views and phased rebuilding.
\begin{semantics}{view-naive}
  \inputleannode{def:view-naive}
\end{semantics}

\begin{semantics}{view-query}
  \inputleannode{def:view-query}
\end{semantics}

\begin{semantics}{view-index}
  \inputleannode{def:view-index}
\end{semantics}

\begin{semantics}{view-indexed}
  \inputleannode{def:view-indexed}
\end{semantics}

\begin{semantics}{fd-syntax}
  \inputleannode{def:fd-syntax}
\end{semantics}

\begin{semantics}{fd-rule}
  \inputleannode{def:fd-rule}
\end{semantics}

\begin{semantics}{fd-instances}
  \inputleannode{def:fd-instances}
\end{semantics}

\begin{semantics}{scheduling}
  \inputleannode{def:phased}
\end{semantics}

% Source proofs: their syntax and checker; the paper's proof section inputs these.
\begin{semantics}[\small]{proofs-syntax}
  \inputleannode{def:proof-syntax}
\end{semantics}

\begin{semantics}{proofs-check}
  \inputleannode{def:proof-facts}
  \inputleannode{def:proof-check}
\end{semantics}

% The provenance pass: instrumentation, the egglog' proofs and their checker, extraction.
\begin{semantics}{provenance}
  \inputleannode{def:instrument}
  \inputleannode{def:log-proofs}
\end{semantics}

% Proof lifting: claims of facts, existence, the query re-reading, and the lift.
\begin{semantics}{lifting}
  \inputleannode{def:lift-claim}
  \inputleannode{def:lift-query}
  \inputleannode{def:lift}
\end{semantics}

\begin{semantics}{lifting-rebuild}
  \inputleannode{def:lift-rebuild}
\end{semantics}
\end{document}
